% =============================================================================
% C-374: Monodromy = 2π 엄밀 증명 구성 및 자명성 판정
% =============================================================================
% 사이클 #374 — 2026-04-27
% 설계자/실행자 (Opus) 산출물
% =============================================================================

\documentclass[11pt]{article}
\usepackage{amsmath,amssymb,amsthm}
\newtheorem{theorem}{Theorem}
\newtheorem{proposition}[theorem]{Proposition}
\newtheorem{corollary}[theorem]{Corollary}
\newtheorem{remark}[theorem]{Remark}
\newtheorem{definition}[theorem]{Definition}

\begin{document}

\section*{C-374: Monodromy Theorem — 엄밀 증명 및 자명성 판정}

% ─────────────────────────────────────────────────────────────────────────────
\section{현황: 기존 논문에 이미 존재}
% ─────────────────────────────────────────────────────────────────────────────

논문 \texttt{unified\_master\_en.tex}에 다음이 이미 포함되어 있음:
\begin{itemize}
  \item \textbf{Theorem (Holonomy Characterization)} [line 2035]:
    반지름 $r$ 원 $\mathcal{C}_r(\rho)$가 단순 영점 $\rho$만 감싸면
    $\mu(\mathcal{C}_r(\rho)) = 2\pi m$ (단, $m$은 영점 차수).
  \item \textbf{Proof} [line 2068]: ``This is a direct application of the
    \textbf{argument principle}.'' 3줄 증명.
  \item \textbf{Remark} [line 2087]: ``Theorem X is the classical argument
    principle restated as a statement about the U(1)-bundle.''
  \item \textbf{Remark (Simple zeros)} [line 2097]: $m=1$ 조건 하에
    $\mu = 2\pi$. 수치 측정 mono/$\pi$ = 2.000 일치.
\end{itemize}

\textbf{결론}: 수학자가 요청한 ``Monodromy = 2$\pi$ 엄밀 증명''은 이미 Paper A에
Theorem으로 수록되어 있으며, argument principle의 직접 귀결임이 명시되어 있음.

% ─────────────────────────────────────────────────────────────────────────────
\section{자명성 판정}
% ─────────────────────────────────────────────────────────────────────────────

\subsection{판정: 반-자명 (Semi-trivial)}

\begin{itemize}
  \item \textbf{수학적 내용}: 자명. Argument principle ($\frac{1}{2\pi i}\oint f'/f\,ds = \#$zeros)
    그 자체의 직접 적용.
  \item \textbf{프레임워크 구성}: 비자명 아님, 그러나 ``의미 있는 재해석.'' $\xi'/\xi$를
    U(1)-다발의 접속으로 해석하고 영점을 위상적 전하(topological charge)로 보는
    \emph{시점 전환}은 Iwaniec--Kowalski, Titchmarsh 어디에도 명시적으로 나타나지 않음.
  \item \textbf{문헌 비교}: 유사한 시점은 Berry (1984, geometric phase), Simon (1983,
    holonomy in quantum mechanics), Witten (1988, topological QFT)에 존재하나,
    $L$-함수 영점 탐지에의 적용은 선행 문헌 없음.
\end{itemize}

\textbf{최종 판정}: Monodromy = $2\pi$ 자체는 \textbf{Theorem}으로 남겨도 됨
(argument principle의 기하학적 재서술). 단, ``새로운 수학적 정리''라고 주장하면 안 됨.
논문에서 이미 ``classical argument principle restated''로 정직하게 기술하고 있으므로
현재 서술이 \textbf{최적}.

% ─────────────────────────────────────────────────────────────────────────────
\section{수학자가 제안한 FE-기반 유도 분석}
% ─────────────────────────────────────────────────────────────────────────────

수학자는 다음 전략을 제안함:
\begin{quote}
  ``$\sigma=1/2$ 위의 적분은 윤곽의 한 변.
  FE 대칭으로 반대편 기여가 동일하므로 합하면 $2\pi$ per zero.''
\end{quote}

이 논증을 엄밀하게 전개하면:

\begin{proposition}[FE-Symmetric Contour Decomposition]
\label{prop:fe_contour}
$\xi(s) = \xi(1-s)$인 완비 $L$-함수에 대해, 단순 영점
$\rho = \tfrac{1}{2} + it_0$을 감싸는 직사각형 윤곽
$\mathcal{R} = [\tfrac{1}{2}-\varepsilon, \tfrac{1}{2}+\varepsilon]
\times [t_0-\delta, t_0+\delta]$ (영점이 $\rho$ 하나뿐이 되도록 $\varepsilon, \delta$ 선택)를
고려하자. 그러면:
\begin{equation}
  \oint_{\mathcal{R}} \frac{\xi'(s)}{\xi(s)}\,ds
  = 2 \int_{\text{right side}} \frac{\xi'(s)}{\xi(s)}\,ds
  + \text{(수평 기여)}.
\end{equation}
함수방정식 $\xi(s) = \xi(1-s)$에 의해, $s = \sigma + it$에서
$\xi'(s)/\xi(s) = -\xi'(1-s)/\xi(1-s)$이므로:
\[
  \int_{\text{left}} \frac{\xi'(s)}{\xi(s)}\,ds
  = \int_{\text{right}} \frac{\xi'(s)}{\xi(s)}\,ds.
\]
따라서 좌변과 우변의 기여가 동일하고, $\varepsilon \to 0$이면 수평 기여가 0으로 수렴하여:
\[
  2\pi i = \oint_{\mathcal{R}} \frac{\xi'(s)}{\xi(s)}\,ds
  = 2 \lim_{\varepsilon \to 0^+}\int_{t_0-\delta}^{t_0+\delta}
  \frac{\xi'(\tfrac{1}{2}+\varepsilon+it)}{\xi(\tfrac{1}{2}+\varepsilon+it)}\,i\,dt.
\]
\end{proposition}

\begin{remark}[이 유도의 한계]
위 분해는 맞지만, 이것이 ``$\sigma=1/2$ 위의 선적분이 $\pi$'' 같은 간단한 결과를
주지는 \textbf{않는다}. 이유:
\begin{enumerate}
  \item 임계선 위($\sigma = 1/2$ 정확히)에서 $\xi'/\xi$는 영점에서 발산하므로
    선적분이 부정적분으로는 정의 불가.
  \item $\xi$가 임계선 위에서 \textbf{실수값}이므로 (FE + 반사 원리),
    $\xi'/\xi$는 순허수. 따라서 $\text{Re}[\xi'/\xi] = 0$ on $\sigma=1/2$.
  \item 영점 사이의 위상 변화는 $0$ 또는 $\pi$ (부호 변환) — 연속적 축적 아님.
\end{enumerate}
결론: FE 대칭은 직사각형 윤곽의 좌/우 대칭을 줄 뿐, ``선적분 = $\pi$''를
의미하지 않음.
\end{remark}

% ─────────────────────────────────────────────────────────────────────────────
\section{비자명 보충 명제: 임계선 제약 (Critical Line Constraint)}
% ─────────────────────────────────────────────────────────────────────────────

기존 Theorem이 cover하지 않는 비자명 사실 하나:

\begin{proposition}[Critical-Line Imaginary Constraint]
\label{prop:imaginary_connection}
$\xi(s)$가 함수방정식 $\xi(s) = \xi(1-s)$ 및 반사 원리 $\overline{\xi(s)} = \xi(\bar{s})$를
만족하면, 임계선 $\sigma = 1/2$ 위에서 (영점 제외):
\begin{equation}
  L\bigl(\tfrac{1}{2} + it\bigr)
  := \frac{\xi'(\tfrac{1}{2}+it)}{\xi(\tfrac{1}{2}+it)}
  \in i\mathbb{R}
  \quad\text{(순허수)}.
  \label{eq:purely_imaginary}
\end{equation}
동치적으로, 접속 $\omega = L\,ds$를 임계선으로 제한하면
$\omega|_{\sigma=1/2} = iA(t)\,dt$ ($A(t) \in \mathbb{R}$).
\end{proposition}

\begin{proof}
$\xi(s) = \xi(1-s)$ 및 $\overline{\xi(s)} = \xi(\bar{s})$로부터,
$s = \tfrac{1}{2} + it$ ($t \in \mathbb{R}$)에서:
\[
  \overline{\xi(\tfrac{1}{2}+it)}
  = \xi\bigl(\overline{\tfrac{1}{2}+it}\bigr)
  = \xi(\tfrac{1}{2} - it)
  = \xi\bigl(1 - (\tfrac{1}{2}+it)\bigr)
  = \xi(\tfrac{1}{2}+it).
\]
따라서 $\xi(\tfrac{1}{2}+it) \in \mathbb{R}$.

$f(t) := \xi(\tfrac{1}{2}+it)$로 놓으면 $f: \mathbb{R} \to \mathbb{R}$.
체인 룰에 의해 $\xi'(\tfrac{1}{2}+it) = f'(t) / i = -if'(t)$. 따라서:
\[
  L(\tfrac{1}{2}+it) = \frac{-if'(t)}{f(t)} = -i\,\frac{f'(t)}{f(t)} \in i\mathbb{R}.
  \qedhere
\]
\end{proof}

\begin{corollary}[이산적 위상 전이]
$\xi(\tfrac{1}{2}+it)$는 실수값이므로, 연속 영점 $t_k < t_{k+1}$ 사이에서
$\arg\xi = 0$ 또는 $\arg\xi = \pi$ (부호에 따라). 단순 영점에서 $\xi$가
부호를 바꾸므로, 위상은 영점마다 $\pi$만큼 불연속적으로 점프한다.
이것이 Berry 위상의 $L$-함수 아날로그이다:
위상 축적은 연속적이지 않고, 영점(위상 전하)에 국소화되어 있다.
\end{corollary}

% ─────────────────────────────────────────────────────────────────────────────
\section{최종 권고}
% ─────────────────────────────────────────────────────────────────────────────

\subsection{논문 내 위치 결정}

\begin{center}
\begin{tabular}{lll}
\hline
내용 & 현재 위치 & 권고 \\
\hline
Monodromy = $2\pi m$ & Theorem (line 2035) & \textbf{유지} (이미 최적) \\
``= argument principle'' & Remark (line 2087) & \textbf{유지} \\
Simple zero $\Rightarrow$ $\mu=2\pi$ & Remark (line 2097) & \textbf{유지} \\
FE contour decomposition & 미수록 & Remark 추가 가능 (선택적) \\
$L|_{\sigma=1/2} \in i\mathbb{R}$ & 미수록 & \textbf{Proposition 추가 권고} \\
이산 위상 점프 & 미수록 & Corollary로 추가 권고 \\
\hline
\end{tabular}
\end{center}

\subsection{비자명성 원천 (무엇이 새로운가)}

\begin{enumerate}
  \item \textbf{프레임워크 구성}: $\xi'/\xi$를 접속으로, $|\xi'/\xi|^2$를 곡률로
    해석하는 것 자체가 $L$-함수 문헌에 선례 없음.
  \item \textbf{정량적 결과}: $\kappa(\rho)/\kappa(\text{non-zero}) \geq 300\times$
    (이것은 Theorem 아닌 Observation이지만 비자명).
  \item \textbf{임계선 제약} (Prop.~\ref{prop:imaginary_connection}):
    FE $\Rightarrow$ 접속이 순허수 $\Rightarrow$ 위상 변화가 이산적.
    이것은 ``왜 영점이 위상적 전하처럼 보이는가''에 대한 기하학적 설명.
  \item \textbf{보편성}: 20+ $L$-함수 (Riemann, Dirichlet, Artin, EC)에서
    동일 기하학 작동. 이것이 Selberg class 통합의 기하학적 근거.
\end{enumerate}

\subsection{수학자에 대한 최종 답변}

\textbf{Q}: Monodromy = $2\pi$가 자명한 재서술인가?\\
\textbf{A}: \textbf{수학적으로 자명} (argument principle). \textbf{개념적으로 비자명}
(다발 기하학 해석은 선례 없음). 현재 논문의 서술이 정확하고 정직함 —
Theorem으로 제시하되 argument principle임을 명시.

\textbf{Q}: 새로 추가할 것이 있는가?\\
\textbf{A}: Proposition~\ref{prop:imaginary_connection} (임계선 위 접속 = 순허수)는
현재 논문에 없으며, FE의 기하학적 의미를 직접 보여주는 비자명 진술.
추가를 권고함. 위치: §3 (Analytic Framework) 또는 §4 (Bundle Theory) 내
Theorem~\ref*{thm:holonomy} 직후.

\end{document}
